% =========================================================================
% PLATFORM TAXONOMY ROUTING (patho.dentnet.space)
% =========================================================================
% DOC_ID: non_odontogenic_cysts
% DOC_SLUG: non_odontogenic_cysts
% DOC_CATEGORY: cat_non_odontogenic_cysts
% DOC_TITLE: الگوریتم جامع: کیست‌های غیرادنتوژنیک و تکاملی فک و دهان
% DOC_ALIASES: non-odontogenic cysts, کیست نازوپالاتین, درموئید, تیرویید زبانی, تیروگلوسال, لنفواپی تلیال, Sistrunk, Dermoid
% =========================================================================
% === CLINICAL EVIDENCE LEDGER & PMIDS ===
% 1. WHO Classification of Head and Neck Tumours (5th Ed, 2024) - Developmental and Non-Odontogenic Cysts.
% 2. Neville BW, Damm DD, Allen CM, Chi AC. Oral and Maxillofacial Pathology (5th Ed).
% 3. PMID: 41721112 - Soares AC et al. (Oral Maxillofac Surg 2026) Clinicopathological Features of Nasopalatine Duct Cyst.
% 4. PMID: 42840700 - Shaaban H et al. (OTO Open 2026) Outcomes of Modified Sistrunk Procedure for Thyroglossal Duct Cyst Removal.
% 5. PMID: 41294898 - Andabak Rogulj A et al. (Diseases 2025) Intraoral Epidermoid and Dermoid Cysts Management.
% 6. PMID: 41562809 - Papadopoulos C et al. (Reports MDPI 2025) Oral Lymphoepithelial Cyst and Branchial Cleft Diagnosis.
% 7. PMID: 37291808 - Kim SJ et al. (Ear Nose Throat J 2025) Endoscopic Transnasal Marsupialization of Large Nasopalatine Cysts.
% =========================================================================

\documentclass[a4paper]{article}
\usepackage[margin=1.15cm,footskip=0pt,headsep=0pt]{geometry}
\usepackage{amsmath,amssymb}
\usepackage{xcolor}
\usepackage{tikz}
\usetikzlibrary{arrows.meta,positioning,shapes.geometric}
\usepackage{hyperref}
\usepackage{fontspec}

\makeatletter
\newcommand{\MathCollectFalse}{}
\newcommand{\MathCollectTrue}{}
\let\@kernel@tabular@init\relax
\makeatother

\usepackage{xepersian}

\settextfont[
  Path = Peyda/,
  BoldFont = Peyda-Bold.ttf,
  ItalicFont = Peyda-Medium.ttf,
  BoldItalicFont = Peyda-Bold.ttf
]{Peyda-Regular.ttf}

\setlatintextfont{Times New Roman}

\IfFileExists{Peyda/Symbola.ttf}{%
  \newfontfamily\symb[Path=Peyda/]{Symbola.ttf}%
}{%
  \IfFileExists{editor/Peyda/Symbola.ttf}{%
    \newfontfamily\symb[Path=editor/Peyda/]{Symbola.ttf}%
  }{%
    \newfontfamily\symb{Symbola}%
  }%
}
\newcommand{\emoji}[1]{{\symb #1}}

\linespread{1.10}
\pagestyle{empty}
\thispagestyle{empty}

\raggedbottom
\setlength{\parindent}{0pt}
\setlength{\parskip}{0pt}

\definecolor{hlgray}{HTML}{E2E8F0}
\definecolor{hlyellow}{HTML}{FEF08A}
\definecolor{hlred}{HTML}{FEE2E2}
\definecolor{darkred}{HTML}{991B1B}
\definecolor{navy}{HTML}{0A2540}
\definecolor{teal}{HTML}{0D9488}

\begin{document}

\begin{center}
{\Large\bfseries فلوچارت تصمیم‌گیری بالینی: کیست‌های غیرادنتوژنیک و تکاملی فک و دهان}\\[2pt]
{\normalsize\bfseries\color{navy} \lr{Non-Odontogenic and Developmental Cysts of the Jaws and Mouth (Differential Triage)}}\\[3pt]
{\small\itshape تریاژ آناتومیک خط وسط، کیست مجرای نازوپالاتین، تیروگلوسال (\lr{Sistrunk})، درموئید کف دهان و لنفواپی‌تلیال}
\end{center}
\vspace{-4pt}
\hrule height 0.8pt
\vspace{6pt}

\begin{center}
\begin{latin}
\linespread{0.95}\selectfont
\begin{tikzpicture}[
    every node/.style={font=\sffamily\footnotesize, align=center},
    stepbox/.style={rectangle, draw=black, fill=white, line width=0.7pt, inner sep=4.5pt},
    decision/.style={diamond, aspect=2.4, draw=black, fill=white, line width=0.7pt, inner sep=2pt, font=\sffamily\scriptsize\bfseries},
    arrow/.style={-{Stealth[scale=0.9]}, thick, black}
]

% 1. Initial Presentation
\node (n1) [stepbox, line width=1.0pt, text width=13.0cm] at (0, 0) {
    \rl{\textbf{تظاهر بالینی و تصویربرداری: تورم کیستیک بافت نرم خط وسط، قدام کام یا لترال گردن}}\\[2pt]
    \scriptsize \rl{رادیولوسنسی قلبی/گلابی‌شکل بین سنترال‌های ماگزیلا $\bullet$ توده خمیری کف دهان یا خط وسط گردن متحرک با بلع}
};

% 2. Mandatory Step 1
\node (n2) [stepbox, text width=14.2cm] at (0, -1.90) {
    \rl{\colorbox{hlyellow}{\strut\textbf{مرحله ۱ الزامی: تصویربرداری سه‌بعدی (\lr{CBCT} / \lr{MRI}) و سنجش حیات پالپ دندان‌های قدامی}}}\\[2pt]
    \scriptsize \rl{رد کیست رادیکولار اینسایزیو با تایید حیات دندان‌های قدامی $\bullet$ بررسی ارتباط توده گردن با استخوان هیوئید یا عضله میلوهیوئید}
};

% 3. Tooth-Related Vitality Decision
\node (d1) [decision] at (0, -4.00) {\rl{دندان بدون حیات؟}\\[-1pt]\scriptsize \lr{Non-Vital Teeth?}};
\node (r_benign) [stepbox, draw=teal, line width=0.9pt, text width=3.8cm] at (5.8, -4.00) {
    \rl{\textbf{کیست پری‌آپیکال رادیکولار \emoji{🔍}}}\\[2pt]
    \scriptsize \rl{منشأ ادنتوژنیک التهابی؛}\\\rl{درمان ریشه (\lr{RCT}) یا}\\\rl{اپیکواکتومی با حفظ پالپ سالم}
};

% 4. Mandatory Surgical Mapping & Anatomic Triage
\node (n3) [stepbox, line width=1.0pt, text width=15.8cm] at (0, -6.40) {
    \rl{\colorbox{hlyellow}{\strut\textbf{تریاژ موضع آناتومیک، آسپیراسیون سوزنی (\lr{FNA}) و مپینگ فضاهای فاسیال \emoji{🔬}}}}\\[2pt]
    \scriptsize \rl{تمایز توده‌های داخل استخوانی ماگزیلا (\lr{NPDC}) از ضایعات بافت نرم ساب‌لینگوال، ساب‌منتال و قوس‌های برانشیال \lr{[\textbf{PMID: 41721112}]}}
};

% 5. Subtype Decision
\node (d2) [decision] at (0, -8.70) {\rl{موقعیت و پاتولوژی}\\[-1pt]\scriptsize \lr{Anatomic Subtype?}};

% 6. 4 Main Columns
\node (h1) [stepbox, text width=3.6cm] at (-6.15, -11.00) {
    \rl{\textbf{مجرای نازوپالاتین}}\\[2pt]
    \scriptsize \lr{Nasopalatine Duct Cyst}\\\rl{رادیولوسنسی متقارن قدام ماگزیلا؛}\\\rl{اپیتلیوم تنفسی/مطبق و اعصاب پالاتین؛}\\\rl{عود با انوکلیشن: \colorbox{hlgray}{\strut\textbf{\rl{کمتر از ۲٪}}}}
};
\node (h2) [stepbox, text width=3.6cm] at (-2.05, -11.00) {
    \rl{\textbf{درموئید / اپیدرموئید}}\\[2pt]
    \scriptsize \lr{Dermoid / Epidermoid}\\\rl{توده قوام خمیری بالای/زیر میلوهیوئید؛}\\\rl{حاوی کراتین، فولیکول مو و سباسه؛}\\\rl{عود موضعی: \textbf{زیر ۳٪}}
};
\node (h3) [stepbox, draw=darkred, line width=1.2pt, text width=3.6cm] at (2.05, -11.00) {
    \rl{\textbf{مجرای تیروگلوسال \emoji{⚡}}}\\[2pt]
    \scriptsize \lr{Thyroglossal Duct Cyst}\\\rl{توده خط وسط گردن متحرک با بلع؛}\\\rl{ریسک بدخیمی کارسینوم پاپیلاری؛}\\\rl{عود بدون سیسترانک:} \colorbox{hlred}{\strut\color{darkred}\textbf{\rl{۵۰٪ تا ۷۰٪}}}
};
\node (h4) [stepbox, text width=3.6cm] at (6.15, -11.00) {
    \rl{\textbf{لنفواپی‌تلیال گردنی/دهانی}}\\[2pt]
    \scriptsize \lr{Lymphoepithelial Cyst}\\\rl{توده زردرنگ کف دهان یا لترال گردن؛}\\\rl{اپیتلیوم سنگفرشی محصور در لنفوئید؛}\\\rl{عود پس از جراحی: \textbf{کمتر از ۱٪}}
};

% 7. Primary Interventions
\node (i1) [stepbox, text width=3.6cm] at (-6.15, -13.50) {
    \rl{\textbf{انوکلیشن از فلپ پالاتال}}\\[2pt]
    \scriptsize \rl{انوکلیشن کامل از طریق دسترسی کام؛}\\\rl{حفظ عصب و عروق نازوپالاتین؛}\\\rl{مارسوپیالیزاسیون در کیست غول‌آسا}\\\rl{\lr{[\textbf{PMID: 37291808}]}}
};
\node (i2) [stepbox, text width=3.6cm] at (-2.05, -13.50) {
    \rl{\textbf{انوکلیشن ترانس‌اورال/گردنی}}\\[2pt]
    \scriptsize \rl{خارج‌سازی با برش داخل دهانی}\\\rl{(در صورت بالای عضله میلوهیوئید) یا}\\\rl{برش اکسترااورال زیر چانه؛}\\\rl{\lr{[\textbf{PMID: 41294898}]}}
};
\node (i3) [stepbox, draw=darkred, line width=1.2pt, text width=3.6cm] at (2.05, -13.50) {
    \rl{\textbf{پروسیجر سیسترانک (\lr{Sistrunk})}}\\[2pt]
    \scriptsize \textbf{\rl{رزکسیون قطعه میانی استخوان هیوئید}}\\\rl{همراه با هسته عضلات قاعده زبان؛}\\\rl{کاهش عود به کمتر از ۳٪}\\\rl{\lr{[\textbf{PMID: 42840700}]}}
};
\node (i4) [stepbox, text width=3.6cm] at (6.15, -13.50) {
    \rl{\textbf{اکسیزیون ساده کپسولی}}\\[2pt]
    \scriptsize \rl{اکسیزیون ترانس‌اورال تمیز در کف دهان؛}\\\rl{یا خارج‌سازی جراحی در حاشیه قدامی}\\\rl{عضله استرنوکلیدوماستوئید}\\\rl{\lr{[\textbf{PMID: 41562809}]}}
};

% 8. Surveillance Regimens
\node (s1) [stepbox, text width=3.6cm] at (-6.15, -15.90) {
    \rl{\textbf{زمان‌بندی پیگیری:}}\\[2pt]
    \rl{\colorbox{hlgray}{\strut\textbf{سالانه تا ۲ سال}}}\\[2pt]
    \scriptsize \rl{رادیوگرافی اکلوزال یا پانورامیک؛}\\\rl{تایید بازسازی استخوان قدام کام}
};
\node (s2) [stepbox, text width=3.6cm] at (-2.05, -15.90) {
    \rl{\textbf{زمان‌بندی پیگیری:}}\\[2pt]
    \rl{\textbf{هر ۶ ماه تا ۲ سال}}\\[2pt]
    \scriptsize \rl{لمس بالینی کف دهان؛}\\\rl{عدم محدودیت حرکتی زبان و تنفس}
};
\node (s3) [stepbox, draw=darkred, line width=1.2pt, text width=3.6cm] at (2.05, -15.90) {
    \rl{\textbf{زمان‌بندی پیگیری:}}\\[2pt]
    \colorbox{hlred}{\strut\color{darkred}\textbf{\rl{هر ۶ ماه تا ۵ سال}}}\\[2pt]
    \scriptsize \rl{سونوگرافی بافت نرم گردن؛}\\\rl{پایش دقیق عود در قاعده زبان}
};
\node (s4) [stepbox, text width=3.6cm] at (6.15, -15.90) {
    \rl{\textbf{زمان‌بندی پیگیری:}}\\[2pt]
    \rl{\textbf{سالانه در معاینات روتین}}\\[2pt]
    \scriptsize \rl{ضایعه خوش‌خیم بدون عود؛}\\\rl{معاینه روتین مخاطی و گردنی}
};

% 9. Clinical Red Flags
\node (redflags) [stepbox, draw=darkred, line width=0.9pt, text width=15.8cm] at (0, -18.25) {
    \rl{\textbf{هشدارهای بالینی و اندیکاسیون‌های ارجاع اورژانسی (\lr{Clinical Red Flags}):}}\\[3pt]
    \scriptsize \colorbox{hlred}{\strut\color{darkred}\textbf{\rl{مداخله فوری در:}}} \rl{(۱) انسداد راه هوایی فوقانی یا دیسفونو/دیسفاژی حاد ناشی از اتساع کیست کف دهان به فضای حلقی $\bullet$ (۲) عدم وجود بافت تیروئید نرمال در بستر آناتومیک اسکن تیروئید قبل از جراحی (\lr{Ectopic Lingual Thyroid}) که جراحی منجر به هیپوتیروئیدیسم دائمی خواهد شد $\bullet$ (۳) رشد سریع توده گردنی تیروگلوسال با لنفادنوپاتی مشکوک به کارسینوم پاپیلاری اولیه مجرا \lr{[\textbf{PMID: 42996571}]}.}
};

% 10. Refractory / Malignancy Protocols
\node (refractory) [stepbox, line width=0.9pt, text width=15.8cm] at (0, -20.40) {
    \rl{\textbf{پروتکل موارد راجعه و پایش بدخیمی نادر مجاری تکاملی (\lr{Malignancy Surveillance}):}}\\[3pt]
    \scriptsize \rl{\textbf{سطح ۱ (کیست عودکننده تیروگلوسال):} ارزیابی مجدد سونوگرافی و \lr{MRI}؛ انجام جراحی مجدد با رزکسیون وسیع‌تر بافت هیویید و بافت همبند قاعده زبان به روش \lr{Extended Sistrunk Procedure}.}\\
    \scriptsize \rl{\textbf{سطح ۲ (پایش بدخیمی کارسینوم پاپیلاری):} پاتولوژی قطعی نمونه جراحی؛ در صورت وجود کارسینوم پاپیلاری تیروئید با تهاجم کپسولی، ارزیابی تیروئیدکتومی توتال و درمان با ید رادیواکتیو (\lr{131-I}).}\\
    \scriptsize \rl{\textbf{سطح ۳ (ضایعات پرفوراتیو نازوپالاتین):} در کیست‌های بزرگ با تخریب وسیع کورتکس بینی، کاربرد تکنیک اندوسکوپیک ترانس‌نازال جهت پیشگیری از فیستول دهانی-بینی.}
};

% Arrows & Routing
\draw [arrow] (n1) -- (n2);
\draw [arrow] (n2) -- (d1);
\draw [arrow] (d1.east) -- node[above, font=\scriptsize\bfseries] {\rl{بله}} (r_benign.west);
\draw [arrow] (d1.south) -- node[right=5pt, font=\scriptsize\bfseries] {\rl{خیر / دندان‌های وایتال}} (n3.north);
\draw [arrow] (n3.south) -- (d2.north);

% Bus-Bar from Decision 2 to 4 Subtype columns
\draw [thick, black] (d2.south) -- (0, -9.80);
\draw [thick, black] (-6.15, -9.80) -- (6.15, -9.80);
\draw [arrow] (-6.15, -9.80) -- (h1.north);
\draw [arrow] (-2.05, -9.80) -- (h2.north);
\draw [arrow, darkred] (2.05, -9.80) -- (h3.north);
\draw [arrow] (6.15, -9.80) -- (h4.north);

% Column Downward Arrows
\draw [arrow] (h1.south) -- (i1.north);
\draw [arrow] (h2.south) -- (i2.north);
\draw [arrow, darkred] (h3.south) -- (i3.north);
\draw [arrow] (h4.south) -- (i4.north);

\draw [arrow] (i1.south) -- (s1.north);
\draw [arrow] (i2.south) -- (s2.north);
\draw [arrow, darkred] (i3.south) -- (s3.north);
\draw [arrow] (i4.south) -- (s4.north);

% Collector Bus into Red Flags Box
\draw [thick, black] (s1.south) -- (-6.15, -17.30);
\draw [thick, black] (s2.south) -- (-2.05, -17.30);
\draw [thick, darkred] (s3.south) -- (2.05, -17.30);
\draw [thick, black] (s4.south) -- (6.15, -17.30);
\draw [thick, black] (-6.15, -17.30) -- (6.15, -17.30);
\draw [arrow] (0, -17.30) -- (redflags.north);

\draw [arrow] (redflags.south) -- (refractory.north);

\end{tikzpicture}
\end{latin}
\end{center}

\end{document}
